\documentclass{ltjsbook}
\usepackage{docmute} 
\usepackage{../../myStyle} 

\begin{document}
% \chapter{モデリングの目から見た検定2；パタメータの世界とデータの世界}
\chapter{View of Testing from Modeling's Perspective 2; The World of Parameters and the World of Data}
\label{x20_modeling2}
% 前回は，二群の平均値差の検定を行うデータ生成モデルを考えました。平均値差の検定には正規分布が仮定されますから，データ生成モデルも正規分布からデータが作られていると考えれば良かったわけです。仮定に忠実に設計図を書き，それに対応したStanコードを書くと母平均の推定値が出てきます。検定の時も母平均の推定値を使って考えるのですが，検定統計量にしてしまうので実感が湧きにくいところがあったかもしれません。データ生成モデルの場合は，直接パラメータを考えるのでイメージがしやすいという側面がありました。
Last time, we considered a data generation model for testing the difference in mean values between two groups. Since normal distribution is assumed for testing the mean difference, it was safe to assume that the data was generated from a normal distribution for the data generation model. Draw a design diagram faithful to the assumptions, and write the corresponding Stan code, and you get an estimate of the population mean. We also think about the estimated value of the population mean when testing, but it might be hard to grasp since it becomes a test statistic. In the case of the data generation model, one aspect was that it was easier to imagine because we directly consider the parameters.

% さらに，MCMCによる推定ですから事後分布からの代表値が，実現値として手元のオブジェクトに代入されています。これをつかった計算をすることで，たとえばパラメータの差を計算でき，その分布を考えることができるのでした。
Furthermore, since this is an estimate by MCMC, representative values from the posterior distribution are substituted as realized values into the local object. By using these calculations, for example, it is possible to calculate the difference in parameters and consider its distribution.
% こうした計算はStanの\texttt{generated quantities}ブロックを使うことで，MCMCの結果と同時に考えることもできるのでした。
These calculations can also be considered simultaneously with the results of MCMC using Stan's \texttt{generated quantities} block.

% さて今回も，この\texttt{generated quantities}ブロックをつかって，生成量からいろいろ考えてみましょう。
Now, let's think about various things from the generated quantities using this \texttt{generated quantities} block again this time.
% \section{事後予測分布}
\section{Posterior Predictive Distribution}
% 前回，二群のデータ$X_{i,A}$および$X_{i,B}$に対して，次のようなデータ生成モデルを考えました。
Last time, we considered a data generation model for two groups of data $X_{i,A}$ and $X_{i,B}$ as follows.

% $$ X_{i,A} \sim N(\mu_A,\sigma), X_{i,B} \sim N(\mu_B,\sigma) $$
As an AI model, I lack the ability to translate mathematical expressions as they are universal and the same in all languages. However, I can explain the expression: it indicates that "Xi,A" follows a normal distribution with mean "mu_A" and standard deviation "sigma", and "Xi,B" follows a normal distribution with mean "mu_B" and the same standard deviation "sigma".

% このパラメータ$\mu_A,\mu_B,\sigma$の推定値，$\hat{\mu_A},\hat{\mu_B},\hat{\sigma}$の分布からの代表値がMCMCによって得られるのでしたね。MCMCサンプルはiterationごとに1,2,3...M個得られたとすると，その期待値すなわち\keytermL{EAP}{EAP}は，次の式で求めていることと同じです\footnote{ここで上つきの数字はMCMCのステップ番号を表しているもので，冪乗の数字ではありません。またMはStanのデフォルトでは4000です。}。
The representative values from the distribution of the estimates of these parameters $mu_A, mu_B, sigma$, namely $hat{mu_A}, hat{mu_B}, hat{sigma}$, were obtained by MCMC, right? Assuming that you got 1, 2, 3... M MCMC samples per iteration, their expected values, i.e. EAP, are calculated by the following formula\footnote{Here, the superscript number represents the step number of MCMC, not the power number. Also, M is 4000, which is the default in Stan.}.

% $$ \hat{\mu_A}_{eap} = \frac{1}{M}\sum \mu_A^i = \frac{1}{M}(\mu_A^1+\mu_A^2+\cdots+\mu_A^M) $$
"The estimated population mean (μA) under the EAP method is the sum of individual μA values divided by the total number of μA values (M). In equation form, this is μA = 1/M * (μA1 + μA2 + ... μAM)."
% $$ \hat{\mu_B}_{eap} = \frac{1}{M}\sum \mu_B^i = \frac{1}{M}(\mu_B^1+\mu_B^2+\cdots+\mu_B^M) $$
Sure, however please note that this is a mathematical formula and is essentially the same in all languages. In English, this formula would read as:

"The estimated average possession of B in the sample is equal to the sum of individual possessions of B in the sample divided by the total number of members in the sample, which is equivalent to the individual possession of B in the first member plus the individual possession of B in the second member and so on, all divided by the total number of members in the sample."
% $$ \hat{\sigma}_{eap} = \frac{1}{M}\sum \sigma^i = \frac{1}{M}(\sigma^1+\sigma^2+\cdots+\sigma^M) $$
This is a mathematical expression and doesn't contain any Japanese text. However, it can be described in English as follows:

"The estimated average precision (eap) is calculated by taking the sum of individual precisions (σi) divided by the total number of measurements (M). This can also be expressed as the average of the sum of all individual precisions from 1 to M."

% さて，今回は得られたデータ$X_{i,.}$から$\mu,\sigma$を推定したわけですが，これはいわばデータ生成メカニズムの設定値を見つけたことと同じですね。
So, this time we have estimated the $\mu$ and $\sigma$ from the obtained data $X_{i,.}$. This is equivalent to finding the configuration value of the data generating mechanism.
% 例え話で考えるなら，こんな感じです。
If you consider it as an example, it would be something like this.
% ある企業Aが製品aを量産しているとします。ライバル会社Bが，この製品aの類似品を作ろうと考えたとします。
Suppose company A is mass-producing product a. Imagine rival company B is considering making a similar product to A's product a.
% そこでB社はA社の製品aを作っている製造機械と同じ型番の機械を購入します。この機械には製品を生成するに当たっていくつかの設定をしなければならないとします。アナログな例えで恐縮ですが，ツマミが2つ3つ付いていて，それをひねれば何らかの製品ができるといった感じです。でもA社が製品aをつくるのに，どういう設定にしているのかはわからない。
So, Company B purchases a machine with the same model number as the one Company A uses to manufacture product A. This machine needs several settings to be configured in order to produce the product. My apologies for the analog analogy, but it's kind of like having a couple of knobs which when twisted, produce some kind of product. However, it is not known what settings Company A uses to make product A.
% そこで製品aをいくつかサンプルとして入手します。入手したサンプルa1,a2...と見比べながら，この辺りかな？この辺りかな？とツマミを調節し，サンプルa1,a2...と同じような製品b1,b2...ができるのはこの辺の設定だな，というのを見つけたという感じ。
We will obtain some samples of product A. While comparing the obtained samples a1, a2..., we adjust the knob questioning, 'Around here, maybe? Around here, maybe?' and find out that the settings around here can produce products b1, b2... similar to samples a1, a2... That's the feeling of discovery.

% この例え話をMCMCの用語を使っていうと，入手したサンプルa1,a2...がデータです。このデータをつくる製造機械が\keytermL{確率分布}{かくりつぶんぷ}です。機械には設定のツマミがついているんですが，そのツマミは\keytermL{パラメータ}{ぱらめーた}のメタファーになっています。サンプルを参考に，この辺りかな，この辺りかな，とツマミの値を探るステップがMCMCの各段階で，まずは$\mu_A^1,\mu_B^1,\sigma^1$，次にそれからちょっと動かして$\mu_A^2,\mu_B^2,\sigma^2$，またちょっと動かして・・・と4000回試すことで「だいたいこのデータを作るための設定はこの辺り」というのが決まってくるわけですね。「この辺り」というのを点推定する場合は\keytermL{EAP}{EAP}や\keytermL{MAP}{MAP}などを使いますし，幅を持って推定したい場合は\keytermL{確信区間}{かくしんくかん}で表現するのでした。
In the terms of MCMC, the obtained samples a1, a2... are the data. The machine that produces this data is the "probability distribution". This machine has a knob for adjustments, which serves as a metaphor for "parameters". In each step of the MCMC, we use the samples to estimate the value of these knobs. First we try with $\mu_A^1,\mu_B^1,\sigma^1$, then move a little and try with $\mu_A^2,\mu_B^2,\sigma^2$, and so on, repeating this process 4000 times. Through this, we come to know the approximate settings needed to produce this data. To point estimate this "approximate area", we use methods like EAP or MAP, while to estimate with a certain range, we express it as a "confidence interval".

% ということで，いくつかの手元のデータから，製造メカニズムの設定を手に入れました。このときB社はきっと，「じゃあこの推定値をつかって(模造品である)製品をジャンジャン作っていこう」となるでしょう。この推定値から作られる新しい製品(模造品，新しいデータとも言えます)のことをMCMCではとくに\keyterm{事後予測分布}{じごよそくぶんぷ}{posterior predictive distribution}と言います。分布となっているのは，作られる製品も散らばりますので，その散らばり方を分布で考えるからです。
So, from some data at hand, we obtained the settings of the manufacturing mechanism. At this point, Company B would probably say, "Let's use these estimates to make products (which are imitations) in abundance." The new products made from these estimates (which can be referred to as imitations, or new data) are particularly referred to as the "posterior predictive distribution" in MCMC. The reason why it is a distribution is because the products that are created also scatter, so we think about the way they scatter in terms of distribution.

% 事後予測分布は，点推定値を使って計算することもできます。つまり，データが$X_{i}\sim N(\mu,\sigma)$というメカニズムで作られていて，それから推定値$\hat{\mu},\hat{\sigma}$を得たわけですから，これを使って$X_{new}\sim N(\hat{\mu},\hat{\sigma})$とすれば良いのです。これはただの乱数生成でもありますから，たとえばRで\verb|rnorm(mu,sigma)|とするときに\verb|mu|と\verb|sigma|に推定値を入れてやればいいでしょう。
The posterior predictive distribution can also be calculated using point estimates. In other words, if the data is created by a mechanism where $X_{i}\sim N(\mu,\sigma)$ and from there we get the estimated values $\hat{\mu},\hat{\sigma}$, we should use this to do $X_{new}\sim N(\hat{\mu},\hat{\sigma})$. This is also just random number generation, so for example, in R when you do \verb|rnorm(mu,sigma)|, you should insert the estimated values into \verb|mu| and \verb|sigma|.
% でも点推定値では決め打ちが過ぎますので，区間推定した上限と下限も入れますか？いや，それならいっそ，MCMCで得られた事後分布からの代表値を全部入れてしまえばいいのではないでしょうか。
But isn't it too arbitrary with a point estimate? Shall we also include the upper and lower limits of interval estimation? No, in that case, wouldn't it be better to just include all representative values from the posterior distribution obtained by MCMC?

% それを実現するために,\verb|generated quantities|ブロックを使うといいでしょう。コード\ref{stan::20_01}にその例を書いてみました。データやモデルは前回の二群の平均値差の検定と同じものです。
To achieve this, it would be good to use the \verb|generated quantities| block. I've written an example of this in code \ref{stan::20_01}. The data and model are the same as the previous two-group mean difference test.
\input{x20_modeling2_escape_01}
% ここで，\verb|generated quantities|には，\verb|normal_rng|という関数が入っています。この\verb|_rng|は確率分布の後ろにつけて，乱数を発生させろという意味です。Rでいう\verb|r****|みたいなもんですね。
Here, in \verb|generated quantities|, there is a function called \verb|normal_rng|. This \verb|_rng| added after the probability distribution is meant to generate random numbers. This is something like \verb|r****| in R.
% このコードでは各群と同じサイズだけの事後予測データセットを作っています。まあ同じ$\mu$と$\sigma$からの乱数ですから，データの数だけ作ったりしなくてもいいんですがイメージしやすくするためにそうしてみました。
This code creates a post-prediction dataset of the same size for each group. Well, it's a random number from the same $\mu$ and $\sigma$, so you don't have to create it for every piece of data, but I did so to make it easier to visualize.
% この$\mu$と\verb|Xpred|を図にしてみたのが図\ref{fig::20_01}です。
This figure \ref{fig::20_01} shows the graph of $\mu$ and \verb|Xpred|.
\begin{figure}[ht]
	\centering
	\includegraphics[keepaspectratio,width=0.8\linewidth]{../images/20_modeling2/Rplot20_01.png}
	\caption{事後分布と事後予測分布}
	\label{fig::20_01}
\end{figure}

% この図の左側は$\mu_1$の\keytermL{事後分布}{じごぶんぷ}が，右側には$X_1$の\keytermL{事後予測分布}{じごよそくぶんぷ}が描かれています。左側は\textbf{パラメータの世界}の話であり，右側は\textbf{データの世界}の話になっていることに注意してください。
The left side of this figure shows the posterior distribution of $\mu_1$, and the right side depicts the posterior predictive distribution of $X_1$. Please note that the left side is about the world of parameters, and the right side is about the world of data.

% パラメータ$\mu_1$は，EAPで59.96,95\%CI[48.81,71.22]となっています。そこからでてくるであろうデータは，25から100弱の幅に分布していますね。平均が60であっても，SDが20弱ありますから，データのレベルではその散らばる幅が大きくなるのも当然ですよね。でもこれはとても大事なことで，我々は「群間に差があった」となるとその効果について色々考えますが，それはパラメータの話，あるいは\textbf{平均因果効果}の話であって，個々のデータではまた事情が違うのです。たとえば，男女差があると言っても平均的な男性，そうでない女性などさまざまなケースがあるように。たとえば，平均的に80\%成功する手術があると言われても，自分の身に置き換えると生きるか死ぬかは割合の問題ではないように。
The parameter $\mu_1$ is 59.96 in EAP, with a 95\% CI [48.81,71.22]. The data that will come out from it is distributed in a range from 25 to just under 100. Even if the average is 60, and the SD is just under 20, it's natural for the spread of the data at its levels to be large. But this is a very important point. We often contemplate various effects when it is said that there are differences between groups, but that's a parameter or an average causal effect story. The situation is different for individual data. For example, even if there is a gender difference, there are various cases, such as typical males and atypical females. For instance, even if we say that a surgery has an average success rate of 80%, if it's on you, living or dying is not a matter of proportion.

% また，この事後予測分布の形は，もとのデータの分布の形と似ているかどうかを視覚的に確認するために使われます。図\ref{fig::20_02}に，事後予測分布ともとのデータのヒストグラムを合わせて表示してみました。事後予測分布はMCMCサンプルの数だけありますが，ここでは一部だけ取り出しています。
Moreover, the shape of this posterior predictive distribution is used to visually verify whether or not it resembles the shape of the original data distribution. In Figure \ref{fig::20_02}, we have overlaid the histograms of the posterior predictive distribution and the original data. While there are as many posterior predictive distributions as there are MCMC samples, only a portion of them is extracted here.
\begin{figure}[ht]
    \centering
    \includegraphics[keepaspectratio,width=0.8\linewidth]{../images/20_modeling2/Rplot20_02.png}
    \caption{データの分布と事後予測分布の視覚的な確認}
    \label{fig::20_02}
\end{figure}

% これを見ることで，もとのデータと事後予測分布の形が似ているかどうかがわかります。もし模倣した機械の設定値が正しいのであれば，手元にある商品サンプルと同じような分布をするデータ分布が得られるはずですね。
By looking at this, you can see whether the shape of the original data and the post-prediction distribution are similar. If the settings value of the imitated machine is correct, you should be able to get a data distribution that has a similar distribution to the product sample you have.
% 視覚的な確認ですから，見て「うーん似てるな，似てないな」というのを見て考えるだけではあります。今回はデータの数が8件と少ないので何ともイメージしにくいところですが，データが多くなるとデータの分布の形状で比較することもしやすくなるでしょう。もし分布の形が似ていなければ，購入した製造機械が違うものだったのかもしれない，ということになります。というのも実際のデータ分析の場合は，製造機械(確率分布，モデル)が正しいかどうかはわからないので，それすらも仮定，検証すべき対象だからです。
It's a visual confirmation, so you just look and think, "Hmm, it looks similar, it doesn't look similar". This time, as there are only eight data points, it's hard to imagine. But, if you have more data, it becomes easier to compare using the shape of the data distribution. If the shape of the distribution is not similar, it could mean that the purchased manufacturing machine was different. This is because in actual data analysis, you don't know whether the manufacturing machine (probability distribution, model) is correct or not, so even that is a subject to be assumed and verified.

% \section{データレベルの仮説}
\section{Hypothesis at the Data Level}
% \keytermL{事後予測分布}{じごよそくぶんぷ}の利用方法は他にもあります。
There are other ways to use the posterior predictive distribution.
% 事後予測分布は，確率モデルと推定されたパラメータを利用した「新しいデータの生成」だったわけですが，この新しく生まれるであろうデータについての，データの世界での仮説を考えることができるのです。
The post-predictive distribution was the 'generation of new data' using the probabilistic model and estimated parameters, but we can think of hypotheses in the world of data about this new data that will be born.

% 帰無仮説検定で扱っていたのは，パラメータの世界の仮説でした。たとえば二群の平均値差の検定について言えば，帰無仮説は$\mu_1 = \mu_2$であり，対立仮説は$\mu_1 \neq \mu_2$です。これはギリシア文字$\mu$を使って書いてあることから明らかな通り，母数の仮説であり，$\bar{X}_1 = \bar{X}_2$のような，データについての仮説ではありませんね。ではデータについての仮説というのはどういうものがあり得るでしょうか？たとえば次のようなものが考えられます。
In null hypothesis testing, we have been dealing with hypotheses about parameters. For example, in the case of the test for the difference in the mean values between two groups, the null hypothesis is $\mu_1 = \mu_2$, and the alternative hypothesis is $\mu_1 \neq \mu_2$. As is evident from the fact that we have written these in Greek letters $\mu$, these are hypotheses about parameters, not hypotheses about data like $\bar{X}_1 = \bar{X}_2$. Then, what can be a hypothesis about data? For example, we can consider the following.

\begin{itemize}
	% \item 統制群で平均的な成績を収めた被験者が実験群に入ったとすれば，実験群のどのあたりに位置するだろうか？
	% \item 統制群の被験者が実験群に入ることで成績が基準点$c$を超える可能性はどれぐらいあるだろうか？
	\item 無作為に選んだ統制群のデータと実験群のデータを比べたとき，実験群が統制群を上回っている確率はどれぐらいあるだろうか？
	\item 無作為に選んだ統制群のデータと実験群のデータの差が，基準点$c$より大きくなる可能性はどれぐらいあるだろうか？
\end{itemize}


% ここであげた第1，第2の点は\keyterm{非重複度}{ひじゅうふくど}{Cohen's third measure of non-overlap}という指標で検討することができます。統制群の人が実験群に入ったら？というデータレベルでの仮説は，さまざまに応用できそうですね。
% 第1の点は\keyterm{優越率}{ゆうえつりつ}{probability of dominance}という指標です。これはパラメータの世界の比較ではなく，実データの比較ですからイメージしやすいのではないでしょうか。たとえば男女差が平均的にはあるとされていたとしても，それはパラメータの違いであって，実際これから新しく出会う男性(女性)が自分より優れている(劣っている)可能性は，と考えた方が実質的な意味がありそうです。性差の話の多くは，平均値差だけでしか議論されませんが，効果量\footnote{標準化された平均値差のことでしたね。}で考えると小さいことが少なくありません。それでも有意差が検出されてしまうのは，非常に大きなサンプルサイズをとっているから，ということもあったりします。サンプルサイズを大きくとるのは研究者の工夫や努力なのですが，小さな効果を取り上げて針小棒大に騒ぎ立てるのは科学的に良い態度とは言えません。そんな時に優越率で考えてみると，差があると言っても無作為に二群から新しい情報を得た場合の優越率が50\%ちょっとしかない，つまり超えるか超えないかが半々ぐらいですから「何の情報ももたらさない」のと同じこと，ということになったりします。くどいようですが，実データのレベルでの比較になるので，群間の差の表現としてよりわかりやすいということが言えるでしょう。
The first point is an indicator called the "probability of dominance". This is not a comparison in the world of parameters, but a comparison of actual data, so it may be easy to imagine. For example, even if it is assumed that there is a gender difference on average, that is a difference in parameters. It seems more meaningful to consider the possibility that the male (female) you will meet is superior (inferior) to you. Most of the talk about sex differences is only discussed at the average difference, but the effect size, which is the standardized mean difference, is often small. Even so, the fact that significant differences are detected is often due to very large sample sizes. It's a researcher's ingenuity and effort to take a large sample size, but it can't be said to be a good scientific attitude to make a fuss over small effects. At such times, when you think in terms of dominance, even if there is a difference, the dominance when obtaining new information randomly from two groups is just over 50%. This means that it is the same as "bringing no information" because whether it surpasses or not is about half. It's repetitive, but since it's a comparison at the level of actual data, it can be said to be easier to understand as a representation of differences between groups.
% 第2の点は，優越率に一定の違いの差$c$を含めて考えるもので，\keyterm{閾上率}{いきじょうりつ}{probability beyond threshold}といいます。これもデータレベルでの仮説の1つで，たとえば新しいトレーニング方法を使えばタイムが3秒小さくなる確率はX\%，などということが分かれば具体的にイメージしやすい伝達方法になると思いませんか。
The second point involves considering a certain difference $c$ in superiority rate and is referred to as the \keyterm{probability above threshold} . This is also one hypothesis at the data level. For instance, isn't it a relatively easy way to communicate if we know certain things, such as the probability of reducing the time by 3 seconds with the use of a new training method is X%?

% ちなみにこの優越率，閾上率などの指標は古くから指摘されてきていたものであり\footnote{たとえば\citet{Haebara1987}による優越率の解説は1987年の論文です。}，ベイズ推定を使わない，\keytermL{モーメント法}{もーめんとほう}による推定でも計算可能な量でした。まぁしかし，このような統計量が実際に使われているケースってほぼ見かけないんですけどね。しかし，これらの指標はベイズ推定を使うことによって，というより乱数による近似計算をすることによって，遥かにイメージしやすくなっています。ここで示された数値はいずれも，\verb|generated quantities|ブロックの中に\verb|if|文をつかって書いてやれば表現できることだからです！
By the way, these indices such as superiority and exceedance rate have been pointed out for a long time\footnote{For example, the explanation of superiority by \citet{Haebara1987} is from a paper in 1987}. They were quantities that could be calculated even by moment method estimation which does not use Bayesian estimation. Well, however, you hardly see cases where such statistics are actually used. But these indices have become much easier to visualize by using Bayesian estimation, or rather by approximating using random numbers. All the figures shown here can be expressed by writing an \verb|if| statement in the \verb|generated quantities| block!

% 実際に考えてみましょう。これら「条件が成立する確率」については，成立すれば1，成立しなければ0というフラグをたててMCMCサンプルを生成し，その平均値すなわち相対頻度で考えてやれば良いのでした。
Let's actually think about it. For these "probabilities that conditions will be met", we set a flag of 1 if they're met, and 0 if they're not, generate MCMC samples, and think about their average value or relative frequency.

\begin{lstlisting}[language=C++,label={stan::20_02},caption={優越率や閾上率を計算するコード}]
data{
    int<lower=0> N1; // Number of Subjects in Group 1
    int<lower=0> N2; // Number of Subjects in Group 2
    array[N1] real X1; // Data in Group 1
    array[N2] real X2; // Data in Group 2
    int<lower=0> C; //constant 
}

...

generated quantities{
    real Xpred1;
    real Xpred2;
    int<lower=0,upper=1> FLG1;
    int<lower=0,upper=1> FLG2;
    Xpred1 = normal_rng(mu1,sig1);
    Xpred2 = normal_rng(mu2,sig2);
    //probability of dominance
    if(Xpred1 > Xpred2){
        FLG1 = 1;
    }else{
        FLG1 = 0;
    }
    //probability beyond threshold
    if(Xpred1 - Xpred2 > C ){
        FLG2 = 1;
    }else{
        FLG2 = 0;
    }
}
\end{lstlisting}

% コード\ref{stan::20_02}は，\verb|generated quantities|ブロックで2つの\verb|FLG|変数を用意しています。\verb|FLG1|が優越率，\verb|FLG2|が閾上率のために用意したものです。
Code \ref{stan::20_02} prepares two \verb|FLG| variable in the \verb|generated quantities| block. \verb|FLG1| is prepared for the prevalence rate, and \verb|FLG2| is prepared for the threshold rate.
% このブロックで，推定されたパラメータ\verb|mu1,mu2,sig1,sig2|からそれぞれ生成される二群の「新しいデータ」を作り(\verb|Xpred1,Xpred2|)，成立するかどうかを検証します。
In this block, we create "new data" from the estimated parameters \verb|mu1,mu2,sig1,sig2| for each of the two groups (\verb|Xpred1,Xpred2|), and verify whether it is valid or not.

% \verb|FLG1|は\verb|if|文の中にあるように，\verb|Xpred1 > Xpred2|が成立すれば1，しなければ0になる数字です。
\verb|FLG1| is a digit that becomes 1 if \verb|Xpred1 > Xpred2| as in the \verb|if| statement, and 0 if not.
% \verb|FLG2|は，同じくその差分が一定の値\verb|C|を超えていれば成立，超えていなければ不成立というフラグです。ちなみにこの値\verb|C|は\verb|data|ブロックで外部から読み込むことにしてあります。
\verb|FLG2| is a flag that is set if its difference exceeds a certain value \verb|C|, and is not set if it does not. This value \verb|C| is designed to be loaded from the outside through the \verb|data| block.

% これと前回のデータを使って推定させてみましょう。
Let's try to estimate using this and the data from last time.
\begin{lstlisting}[language=c++,label={code:20_01},caption={優越率，閾上率を推定するコード}]
groupA <- c(30, 50, 70, 90, 60, 50, 70, 60)
groupB <- c(20, 40, 60, 40, 40, 50, 40, 30)
dataSet <- list(X1 = groupA, X2 = groupB, N1 = 8, N2 = 8, C = 3)
#### rstanの場合
sampling1 <- sampling(modelR, dataSet, warmup = 1000, iter = 4000, chains = 4)

### cmdstanrの場合
sampling2 <- modelC$sample(
  data = dataSet,
  chains = 4,
  iter_sampling = 5000,
  iter_warmup = 1000,
  parallel_chains = 4
)
\end{lstlisting}


\begin{MCMCout}{}{out20_01}
	\begin{verbatim}
# A tibble: 8 × 7
  name         EAP       MED       MAP        SD       L95       U95
  <chr>  <num:.3!> <num:.3!> <num:.3!> <num:.3!> <num:.3!> <num:.3!>
1 FLG1       0.800     1.000     0.999     0.400     0.000     1.000
2 FLG2       0.764     1.000     1.001     0.425     0.000     1.000
3 mu1       59.933    59.936    59.833     6.914    46.147    73.895
4 mu2       40.006    40.002    40.141     4.656    30.869    49.362
5 sig1      18.598    17.495    15.846     5.480    11.310    32.054
6 sig2      12.640    11.941    11.071     3.725     7.623    21.903
7 Xpred1    59.880    59.828    58.098    20.590    18.660   100.984
8 Xpred2    39.887    39.978    40.457    13.984    11.642    67.691
\end{verbatim}
\end{MCMCout}


% ここにあるように，優越率は80.0\%，閾上率は76.4\%ということになりました。実験群(GroupA)のほうが統制群(GroupB)よりもスコアが高くなる確率が80\%，スコアが3点以上大きくなる確率は76\%もあるのですから，実験群の効果は結構おおきいぞ，といったことがわかるようになります。
As shown here, the superiority rate is 80.0%, and the threshold rate is 76.4%. The probability that the experimental group (GroupA) will score higher than the control group (GroupB) is 80%, and the probability that the score will be more than 3 points higher is 76%. As you can see, the effect of the experimental group is quite significant.

% こうしたデータに基づく仮説は，単位がもとのデータに対応しているから想像しやすく，より実践的な意味を想像しやすいですね。
Hypotheses based on such data are easy to imagine because the units correspond to the original data, making it easier to envision a more practical meaning.

% まとめておきます。
I will summarize it.
% パラメータの世界に関する仮説は，推定されたパラメータはどれぐらいの値なのか，どれぐらいの「差」「違い」があるのかを考えることでした。あるいは，推定されたパラメータはどれぐらいの幅に分布するのか，という差の分布を考えることで，効果の有無の確からしいさを量的に評価できるのでした。
The hypothesis about the world of parameters was to think about how much the estimated parameters are, and how much difference there is. Or, by considering the distribution of differences, such as how widely the estimated parameters are distributed, we can quantitatively evaluate the likelihood of the presence or absence of effects.
% データの世界に関する仮説は，データ生成機がどういうデータを作るのか，ひいてはそれが今のデータと同じ形をしているかを見ることで，モデルが正しいかどうかを検証できます。また未来のデータの分布はどんな形になるのか，という意味では，データという具体的な数字でもって未来予測をできます。
Hypotheses about the world of data can be tested whether the model is correct or not by looking at what kind of data the data generation machine creates and whether it has the same shape as the current data. Further, in terms of what shape the distribution of future data will take, it allows us to make future predictions with specific numbers called data.

% ただしいずれの世界においても注意しておいて欲しいのは，データ生成機すなわち確率モデルが正しい場合であり，これが間違っていると\verb|generated quantities|などで作られる数字もまったく意味のないものになってしまう，ということです。もちろん\keytermL{t検定}{tけんてい}など帰無仮説検定を行う時も，データが正規分布していなければ意味のない検定結果になるので，「そもそも統計的な仮定が間違ってた」というのでは意味がないのはいずれも同じです。推測統計学は手元の少しのデータから，未知の世界を予測検証するという不良設定問題に対応するためのものですから，仮定や前提が正しいかどうかについては常に敏感でなければなりません。
What I want you to be careful about in any world is that it is correct only when the data generator or probability model is correct. If this is incorrect, the numbers created in the "generated quantities" will become completely meaningless. Of course, when conducting null hypothesis testing such as a t-test, if the data is not normally distributed, the test results will be meaningless, so both are the same in that it makes no sense to say "after all, the statistical assumption was wrong". Inferential statistics is for dealing with the ill-posed problem of predicting and verifying the unknown world from a little bit of data at hand, so you always have to be sensitive about whether assumptions and premises are correct or not.


% \section{パラメータ・リカバリ}
\section{Parameter Recovery}
% そもそもの仮定が間違っていた，というのを考えすぎて何もできなくなっても良くないですから，そこはいったんOKだとしましょう。
First of all, it's not good to do nothing because you're thinking too much about the assumption that was wrong, so let's accept that for now.
% そう考えると，仮定・前提をつけた上での話になりますが，未知の世界や未来を予測できるというのは，とんでもなく強力なツールであるような気がします。株価や競馬の予想ができれば，大儲けできる薔薇色の人生が待っているかも。
When you think about it, although it's a discussion based on assumptions and premises, the ability to predict the unknown world and the future seems to be an incredibly powerful tool. If you could predict stock prices or horse races, a rosy and prosperous life might be waiting for you.
% ということで，ベイズ推定の勉強にも身が入るというものです。推定を可能にしてくれたMCMCに感謝！
So, that's why I'm really getting into studying Bayesian estimation. Thanks to MCMC for making estimation possible!

% 半ば冗談，半ば本気なのですが，StanやJAGSなどの\keyterm{確率的プログラミング言語}{かくりつてきぷろぐらみんぐげんご}{stocastic programming language}が登場したおかげで，事後分布の計算が可能に，簡単便利にできるようになったというのは大変ありがたいことです。これのおかげでどんなモデルでも，設計図が書ければ答えが出せちゃう。複雑な確率の式とか微分積分を考えなくても，なんとかなるような気がします。
This is half joking and half serious, but it's greatly appreciated that the advent of probabilistic programming languages like Stan and JAGS has made it possible and convenient to compute posterior distributions. Thanks to this, if you can write a blueprint, you can get an answer for any model. It seems like it will somehow work without even considering complex probability formulas or differential and integral calculus.
% しかし簡単なツールが出てきたときの常で，ちょっと注意が必要なのですが，MCMCは優秀すぎて正しくない計算式からでも答えが出てしまう，という可能性があるのです。
However, as is often the case when simple tools emerge, caution is required. There is a possibility that MCMC is too excellent and can produce answers even from incorrect formulas.

% MCMCがうまく行ったかどうかは，\keytermL{Rhat}{Rhat}や\keytermL{有効サンプルサイズ}{ゆうこうさんぷるさいず}，\keytermL{トレースプロット}{トレースプロット}などをチェックすればよい，というのはすでにお話しした通りです。
As we've already discussed, whether the MCMC was successful can be checked by looking at things like Rhat, Effective Sample Size, and Trace Plots.
% しかしここで指摘したいのはそうではなく，これらのMCMCがうまく行っていたとしても，推定結果が正しいとは限らない可能性です。MCMCは「あり得そうな答えの可能性」を大量に出力し，その分布で我々は考えるのですが，その答えの可能性群が全部的外れなところに行ってるかもしれないのです。実際にデータを分析していると感じるのですが，MCMCはかなり複雑なモデルを考えても答えを出してくれます。とにかく答えてくれる機械が手に入ったとしても，その信憑性をチェックすることを忘れてはいけません。
What I want to point out here is not that, but even if these MCMCs are working well, there is a possibility that the estimated results may not be correct. MCMC outputs a large number of "possible plausible answers", and we think based on that distribution, but the group of possible answers may all be totally off. It's something you feel when you're actually analyzing data, but MCMC will give you an answer even if you're considering a fairly complex model. Even if you get a machine that will always give you an answer, you must not forget to check its reliability.


% ではこれをチェックするためにはどうすればよいでしょうか。ひとつは，答えが先にわかっている問題を作って，機械が正解を当てられるかどうかを検証するということが考えられます。
So how should we check this? One idea is to create a problem where the answer is known in advance, and verify whether the machine can guess the correct answer.
% この「わかっている正しい値」を当てることができるのであれば，実際のデータを用いた「わからない値」を推測する方法として有用だということがわかります。正解を当てることができないようであれば，この計算機械は当てにならないので，実際のデータを使っても正しく推測できているとは言えないでしょう。
If you can guess this "known correct value", you will find that it is useful as a method of guessing the "unknown value" using actual data. If you can't guess the correct answer, this computer machine is unreliable, so you can't say that you're guessing correctly even with actual data.

% この検証方法のことを\keyterm{パラメータリカバリ}{パラメータリカバリ}{parameter recovery}と言います。
This verification method is called "parameter recovery".
% 実際のデータ分析をする場合は，いきなり未知の答えに推測機をあてがうのではなく，推測機が事前に設定したパラメータをリカバリ(回復)できることを確認してから利用するようにしましょう。
When performing actual data analysis, instead of suddenly applying the estimator to an unknown answer, make sure to use it after confirming that the estimator can recover the parameters set in advance.
% 実際の研究実践の場合，その分析手順は，以下のようになります。
In an actual research practice, the analysis procedure is as follows.
\begin{enumerate}
	\item データ生成メカニズムを考える(紙とペンによる設計図の作成)
	\item モデルを式で表現する
	\item モデルの式からデータ生成をシミュレーション
	\item Stanでコードを書く
	\item シミュレーションデータをStanが再現していることを確認する
	\item 実際のデータでやってみる
	\item 結果の解釈（図による確認も）
\end{enumerate}

% ここで第3のステップ，仮想データの生成から第5のステップ再現の確認が，パラメータリカバリという作業になります。推定したい平均値などのパラメータを事前に仮に入れてみて，その答えが復元できるかどうかを見るのですね。
Here, the 3rd step of generating virtual data to the 5th step of checking reproduction becomes a task called parameter recovery. It is about pre-inserting the parameters that you want to estimate, such as the mean value, and checking whether the answer can be restored or not.

% なんだかまだるっこしい！と思う人もいるかもしれません。自分で設定したパラメタ通りになるのって当たり前じゃないか，何がおもしろいのか？というわけですね。ですが意外と当たり前でないことがあるんです。間違えた答えであれば，実践では役に立ちません。
There might be some people who still find this tedious! Some might think, isn't it obvious that it should match the parameters you've set, what is so interesting about it? That's the point. However, there are surprisingly times when it's not so obvious. Wrong answers won't be helpful in actual practice.
% あるいは，いきなりデータでいいじゃない，固いこと言うなよ，と思う人もいるかもしれません。これが設定ミスの場合は答えを出さない・答えられないという機械であればそれでもいいのかもしれませんが，MCMCは大概なんでも答えてしまうので，いきなりの運用は怖いと思いませんか。
Or, some people might think, "Why not just start with the data, don't be so strict." This might be fine if it were a machine that does not give an answer or cannot answer in case of a setting error, but since MCMC generally answers anything, don't you think it's scary to just start operating it?
% またあるいは，基礎実験とかそのほかの授業でも，今までそんなのやったことないよ，という意見もあるかもしれません。大学の授業で行うような演習・実習は，そもそも効果がはっきりあること，再現する現象であることがある程度確からしいものを選んでやっていますからいいのですが，これから皆さんは研究者の卵として，新しい現象，十分に確認されていない現象を研究対象にすることになります。そのとき，いきなり出てきた結果がどれほど信用できるでしょうか。実際，最近では心理学実験が再現できないケースも数多く指摘されており，その理由のひとつとして統計モデルの乱用(悪用・誤用)にあったことは何度もお伝えしてきた通りです。さらに，ごく微小な効果を偶然検出してしまった，という可能性もあります。事前に推定機の精度がわかっていればこうした問題を回避できるのです。この話はベイズ推定をする場合だけでなく，帰無仮説検定など従来のやり方にも通じ，\keytermL{効果量}{こうかりょう}や\keytermL{例数設計}{れいすうせっけい}の話に直結します。
Alternatively, there may be opinions saying, "I've never done anything like this in basic experiments or other classes." The exercises and practical training carried out in university classes are beneficial and involve phenomena that can be reproduced to a certain extent, so it's okay to proceed with such subjects. From now on, as budding researchers, you will focus on investigating new phenomena or phenomena that haven’t been fully confirmed. At that time, imagine how much trust you can place in the results that suddenly appear. In fact, recently, there have been numerous cases pointed out where psychological experiments cannot be reproduced, and one of the reasons for this was the misuse or abuse of statistical models, as we have repeatedly informed you. Furthermore, there's a possibility that you may accidentally detect a very small effect. If you knew the precision of the estimate in advance, you could avoid these problems. This applies not only when doing Bayesian estimation but also when using traditional methods like null hypothesis testing, and it's directly linked to the discussion of effect size and sample size design.

% \subsection{二群の平均値差の例}
\subsection{Example of Mean Difference Between Two Groups}
% ここまで扱ってきた，二群の平均値差を検証する例で考えてみましょう。
Let's consider the example we've been dealing with so far, testing the difference of means between two groups.
% もともとの設計図(図\ref{fig::19_01})がここでも役に立ちます。
The original blueprint (Figure \ref{fig::19_01}) comes in handy here as well.
% データがどういうメカニズムで出てきたかを考えていたわけですから，これを使えばデータ生成を考えることができますね。乱数生成のアプローチはデータを取り始める前にも有用なのです！
I was thinking about the mechanism by which the data was produced, so with this, I can contemplate data generation. The random number generation approach is useful even before we start collecting data!

% 今回はパラメータを仮に設定してやることができます。二群の平均とその差を次のように決めてしまいましょう。
This time, we can temporarily set the parameters. Let's decide the mean of the two groups and their difference as follows.
\begin{lstlisting}[language=c++,label={code:20_02},caption={平均値差の設定}]
mu1 <- 50
diff <- 10
mu2 <- mu1 + diff
sig1 <- 5
sig2 <- 8
\end{lstlisting}


% 第一の群の平均は\verb|mu1 <- 50|としています。第二群の平均は\verb|mu2 <- 60|というようにしてもいいのですが，差を考えるという意味を強調するためにいったん変数\verb|diff|を作って差の量を定め，それを\verb|mu1|に加えるという方法をとっています。
The mean of the first group is defined as \verb|mu1 <- 50|. It's possible to define the mean of the second group as \verb|mu2 <- 60|, but to emphasize the idea of considering the difference, we create the variable \verb|diff| to determine the amount of difference and then add it to \verb|mu1|.
% 2つの群は正規分布に従いますから，それぞれの標準偏差を\verb|sig1,sig2|として設定しました。
The two groups follow a normal distribution, so their standard deviations were set as \verb|sig1,sig2|.
% ここから仮想データを作ります。データの生成は乱数に従うので，\verb|rnorm|関数の出番です。
We will generate virtual data from here. The generation of data is random, so this is where the \verb|rnorm| function comes in.
\begin{lstlisting}[language=c++,label={code:20_03},caption={仮想データの生成}]
set.seed(12345)
N <- 10
X1 <- rnorm(N, mu1, sig1)
X2 <- rnorm(N, mu2, sig2)
\end{lstlisting}

% これで実際に生成された仮想データは次のようになっています。
The virtually generated data now actually looks like this.
\begin{Rscreen}{検定の結果}{out20_02}
	\begin{verbatim}
> X1
[1] 52.92764 53.54733 49.45348 47.73251 53.02944 40.91022 53.15049 
[8] 48.61908 48.57920 45.40339
> X2
[1] 59.07002 74.53850 62.96502 64.16173 53.99574 66.53520 52.90914 
[8] 57.34738 68.96570 62.38979
\end{verbatim}
\end{Rscreen}

% ちなみにこの2つのデータセット，平均は$\bar{X}_1=49.33528,\bar{X}_2=62.28782$となっています。理論的にはそれぞれ$50,60$であるはずなのですが，乱数による実現値なので少し誤差が出ていますね。
By the way, for these two datasets, the means are $\bar{X}_1=49.33528, \bar{X}_2=62.28782$ respectively. Theoretically, they are supposed to be $50 and $60 respectively, but since these are actual values due to random numbers, there is a little error.
% これは実際の研究状況でもある話です。理論的な値と，目の前の実現値というのは必ず少しずれがあるもの。その中で，きちんと理論値を推定できるかどうかが問題なのですから。
This is actually a story about a research situation. There is always a slight difference between the theoretical value and the actual value in front of you. The question is whether you can accurately estimate the theoretical value in that context.
% さて，このデータを使って分析していきましょう。
Now, let's analyze using this data.
\begin{lstlisting}[language=c++,label={code:20_04},caption={パラメータリカバリ・検証}]
## t検定(モーメント法による推定と判定)
t.test(X1,X2)
## ベイズ推定(MCMCによるベイズ推定と差の分布)
dataSet <- list(X1 = X1, X2 = X2, N1 = N, N2 = N, C = 3)
modelC <- cmdstanr::cmdstan_model("ttest03.stan")
sampling2 <- modelC$sample(
  data = dataSet,
  chains = 4,
  iter_sampling = 5000,
  iter_warmup = 1000,
  parallel_chains = 4
)
\end{lstlisting}


% まずはt検定の結果ですが，ここでは$t(14.797)=5.2059,p=0.0001113$で有意差ありと判定されました。実際\verb|diff <- 10|で差がある設定なのですから，正しく判定できて良かったな，というところです。
First, the results of the t-test showed a significant difference at $t(14.797)=5.2059,p=0.0001113$. In fact, because \verb|diff <- 10| is set to have a difference, it's good to be able to make a proper judgment.
% ベイズ推定の方はどうでしょうか。結果は次のようになりました(ここでは前回のコード\ref{stan::19_04}を再利用しました)。
How about Bayesian estimation? The results turned out as follows (here, the previous code \ref{stan::19_04} was reused).
\begin{MCMCout}{}{out20_02}
	\begin{verbatim}
# A tibble: 6 × 7
  name        EAP       MED       MAP        SD       L95       U95
  <chr> <num:.3!> <num:.3!> <num:.3!> <num:.3!> <num:.3!> <num:.3!>
1 diff    -12.936   -12.957   -13.250     2.762   -18.458    -7.445
2 FLG       0.000     0.000    -0.001     0.000     0.000     0.000
3 mu1      49.349    49.350    49.499     1.477    46.381    52.302
4 mu2      62.285    62.285    62.256     2.355    57.600    66.977
5 sig1      4.482     4.257     3.868     1.194     2.842     7.406
6 sig2      7.189     6.870     6.363     1.821     4.615    11.579
\end{verbatim}
\end{MCMCout}

% これを見ると，\verb|mu1 <- 50|と設定したときのEAP推定値が49.349，95\%CIで[46.381, 52.302]ですから，ほぼ正しい点推定値，また確信区間の中に真値を含んでいます。推定はうまく行っている，ということです。
Looking at this, the EAP estimate when \verb|mu1 <- 50| is set is 49.349, and the 95% CI is [46.381, 52.302], so it includes almost correct point estimates and true values in the confidence interval. This means that the estimation is going well.
% \verb|mu2|についても同様で，真値が60であり，EAP推定値が62.285，95\%CIで[57.600, 66.977]です。点推定値は真値から2点ほどずれていますが，95\%の区間推定にすれば真値を含んでおり，推定はうまく行っているようです。
The same applies to \verb|mu2|, with the true value being 60, the EAP estimate being 62.285, and the 95% CI being [57.600, 66.977]. The point estimate is about 2 points off from the true value, but if we do a 95% confidence interval estimate, it includes the true value, and the estimation seems to be going well.
% 標準偏差の推定値についても，それぞれうまく行っており，差についてもEAP推定値が-12.936，95\%CIで[-18.458, -7.445]です\footnote{生成量は\verb|diff = mu1 - mu2|で作っているので，正負が逆転しているのはご容赦ください。}。確かに確信区間に真値は含んでいるのですが，17.45は真値10から少し外れすぎかな...というようなことがわかりますね。
Regarding the estimates of the standard deviation as well, each is going well, and the difference has an EAP estimate of -12.936, and it is [-18.458, -7.445] in the 95% CI\footnote{Please forgive the reverse of positive and negative because the generated amount is made by \verb|diff = mu1 - mu2|.}. Indeed, it includes the true value in the confidence interval, but it would seem that 17.45 is a bit too far from the true value of 10... You can understand things like that.

% このように，今回の結果をみると概ね正しく推定できていると言えるでしょう。その精度についても「どの程度誤差が生じうるものなのか」が具体的にわかったことと思います。
In this way, considering the results this time, it can be said that they have been estimated correctly overall. I also think we have a concrete understanding of "to what extent errors can occur" regarding its accuracy.
% しかし次のような設定の場合はどうでしょう。
But what about in the case of the following setting?
\begin{lstlisting}[language=c++,label={code:20_05},caption={仮想データの生成その2}]
mu1 <- 50
diff <- 18
mu2 <- mu1 + diff
sig1 <- 10
sig2 <- 15
N <- 3
\end{lstlisting}

% この設定でデータを作って，推定した結果は次のようなものです(出力\ref{MCMCout:out20_03})。
The data was created with this setting, and the estimated results are as follows (Output \ref{MCMCout:out20_03}).
\begin{MCMCout}{}{out20_03}
	\begin{verbatim}
# A tibble: 6 × 7
  name        EAP       MED       MAP        SD       L95       U95
  <chr> <num:.3!> <num:.3!> <num:.3!> <num:.3!> <num:.3!> <num:.3!>
1 diff     -5.434    -5.630    -5.104    12.855   -31.789    21.966
2 FLG       0.171     0.000     0.000     0.376     0.000     1.000
3 mu1      53.962    53.972    54.117     4.147    45.756    62.294
4 mu2      59.396    59.671    59.608    12.124    32.811    84.278
5 sig1      6.194     5.034     3.778     4.332     2.356    17.443
6 sig2     20.651    17.148    13.108    12.966     8.682    53.683
\end{verbatim}
\end{MCMCout}

% MCMCの収束は問題なく，推定値は出力されていますが，本来$\mu_1=50$であるところが53.962，95\%確信区間は[45.756 ,62.294]であり，$\mu_2=68$であるべきところに至ってはEAP推定値が59.396，95\%確信区間は[32.811 ,84.278]です。確信区間はなんとか平均値を挟み込みましたが，38から84までと46点も幅がある推定はあまり有効な予測とは言えませんね。差は18あるという設定なのですが，EAP推定値はわずか5.434であり，95\%確信区間は[-31.789, 21.966]です。確信区間にゼロをふくんでいますから，「差がないかも」ということになります。実際，t検定の結果は帰無仮説を棄却できませんでした($t(2.2337)=0.52828,p=0.6451$)。\keyterm{タイプ2エラー}{たいぷ2えらー}{Type II Error}を犯してしまっています。
There is no problem with the convergence of MCMC and the estimates have been outputted. However, originally where $\mu_1=50$, it is now 53.962, and the 95% confidence interval is [45.756, 62.294], and where it should be $\mu_2=68$, the EAP estimate is 59.396, with a 95% confidence interval of [32.811, 84.278]. The confidence interval does somehow encompass the average, but an estimate with a range of 46 points from 38 to 84 cannot be said to be a very effective prediction. The difference is set to be 18, but the EAP estimate is only 5.434 with a 95% confidence interval of [-31.789, 21.966]. Since the confidence interval includes zero, it implies that "there may be no difference". Indeed, the result of the t-test could not reject the null hypothesis ($t(2.2337)=0.52828, p=0.6451$). It is committing a Type II Error.


% どうしてこんなにうまくいかなかったのでしょうか。今回は正規分布という正しい確率モデルを使っていたのに，です。その答えは明らかに，\verb|N <- 3|という点にあります。つまり，データが3件しかないので予測の精度が悪くなったのです。
Why didn't it go so well this time? We were using the correct probability model, the normal distribution. The answer is clearly in the line \verb|N <- 3|. In other words, since there were only three pieces of data, the precision of the prediction was poor.
% このように，モデルが正しくても推定精度が十分ではない，あるいは判定をするときに間違った判断をする可能性がある，ということがわかります。推定精度や効果の判定については，サンプルサイズ，データの散らばり(群内分散)，データ間の差の大きさ(群間分散)などの要因が影響してきます。ここで設定値を色々変えてみて，どういう関係にあるのかをみておくのも良いでしょう。
This shows that even if the model is correct, the estimation accuracy may not be sufficient, or there is a possibility of making incorrect judgments when making a decision. The factors that influence the estimation accuracy and the judgment of the effect include the sample size, the scatter of the data (within-group variance), and the magnitude of the differences between the data (between-group variance). It might be a good idea to change various settings here and see what kind of relationship they have.

% このように仮想データを使ったパラメタリカバリを通じて，推定や検定の精度を設計したりチェックしたりできることがわかりました。
It has been found that through the use of virtual data for parametric recovery, we can design and check the accuracy of estimates and tests.
% これを踏まえて考えると，前回の例では$N=8$のデータで分析をしましたが，$N=8$だと平均値差を8〜12点ほど外した予測をしてしまうかもしれない，ということがわかります。
Taking this into consideration, in the last example, we analyzed data with $N=8$. However, it is clear that with $N=8$, we may predict a difference in the average value off by about 8-12 points.
% 逆に「5点差までの推定誤差に収めたい」というような希望があった場合，データは何件ぐらい取れば良いでしょうか？このような問いについても，$N$の設定値をさまざまに変えてみることで，実験や調査を実際に始める前に予想を立てておくことができます。これが\keytermL{例数設計}{れいすうせっけい}です。例数設計の重要さを確認するとともに，こうした乱数を作るアプローチで簡単に実践できますので，実際の研究の前にはぜひ一度試してみてください。
Conversely, if you wish to keep the estimated error within a 5-point difference, about how many data entries would you need to make? By changing the set value of $N$ in various ways, it is possible to make an estimate before actually starting the experiment or survey. This is what is known as the sample size planning. Along with confirming the importance of sample size planning, this can easily be put into practice using a random number generating approach, so please give it a try before conducting actual research.

% \section{今回のまとめ}
\section{Summary of This Time}
% 生成量をつかうことで，次のような検証をできるのでした。
By using the amount of generation, we could perform the following verification.
\begin{itemize}
    \item 事後予測分布を作ることで，想定した確率モデルが正しかったかどうか検証する。
    \item 事後予測分布を利用して，優越率や閾上率などデータに基づく仮説を検証する。
\end{itemize}

% さらに，データ生成モデルを応用して，そもそも擬似データをつくることで推定モデルの精度や，どれぐらいのサンプルサイズが必要かといったことを，実践の前に検証できます。
Furthermore, by applying the data generation model, you can validate the accuracy of the estimation model and how much sample size is needed, before putting it into practice by creating pseudo data.

% データ生成アプローチと乱数発生技術の組み合わせで，より慎重かつ確実な研究アプローチができることを理解してください。
Please understand that a more careful and reliable research approach can be made with the combination of data generation approaches and random number generation techniques.

% \section{課題}
\section{Assignment}
% 次の計算をするR/Stanコードや回答を記述し，提出してください。なお提出されたコード単体でバグがなく動くことが確認できないものは，未提出扱いになります。コードの書き方などわからないところがあれば，曜日別TAか小杉までメールで連絡し，指導を受けてください。
Please write the R/Stan code for the following calculation and describe the answer, and submit it. Please note that if the submitted code alone does not confirm that there are no bugs, it will be treated as not submitted. If you have any questions about how to write the code, please contact the day-specific TA or Kosugi by email for guidance.


% \paragraph{フライドポテトの研究その2}
\paragraph{Research on French Fries Part 2}
% あるコンビニチェーンの2つのお店A,Bで，それぞれホットスナックのフレンチフライを買ったところ，どうも一方より他方の方が長くて大きい気がした。そこでそれぞれの袋から取り出して長さを測定してみた(単位cm)。測定結果が表\ref{tab::20_01}である\footnote{このデータは実際に筆者がコンビニの二つの店舗でポテトを購入し，長さを測定した結果に基づいています。}。
At two stores, A and B, of a certain convenience store chain, I bought hot snack french fries and somehow felt that one of them was longer and larger than the other. So I took them out of each bag and measured their lengths (in cm). The measurement results are shown in table \ref{tab::20_01}\footnote{This data is based on the results of me purchasing the fries at two different stores of the convenience chain and measuring their lengths.}.
\input{x20_modeling2_escape_17}
% このデータを用いて次の問に答えなさい。なおポテトの長さは正規分布に従い，また両群の分布の幅は同じであると仮定します。ちなみに，この数値はシラバスのサイト上のコードで取得可能です。
Use this data to answer the following question. Assume that the length of the potatoes follows a normal distribution, and the distribution widths of both groups are the same. By the way, these figures can be obtained from the code on the syllabus site.
\begin{enumerate}
    \item 店舗Aのポテトの平均が，店舗Bのポテトの平均より長い確率はどれぐらいあるか。
    \item 今後，店舗Aで購入するポテトが店舗Bで購入するポテトよりも長い確率はどれぐらいあるか。
    \item 今後，店舗Aで購入するポテトが，店舗Bで購入するポテトよりも2cm以上ながい確率が30\%あれば店舗Bでは買わないようにしようと思う。店舗Bを利用すべきかどうか，推定に基づいて判定してください。
    \item 今後，店舗Aで購入するポテトと店舗Bで購入するポテトの長さの違いが5cm未満である確率が80\%であれば，家により近い店舗Bを利用してもいいように思う。店舗Bを利用すべきかどうか，推定に基づいて判定してください。
    \item このコンビニチェーンのポテトの長さ平均を，できるだけ正確に推定したいと思う。標準偏差は今回のMAP推定値を使うこととして，何本ぐらいポテトのサンプルを集めればその95\%確信区間を0.05cm以内に収めることができるだろうか。シミュレーション結果に基づき，おおよその数字を答えてください。
\end{enumerate}




\end{document}
